\section{Matched readouts and complete recovery}\label{sec:training69}
For $x\in[0,1]^d$, the two fixed feature maps used in the comparison are
\[
 \phi^{\rm R}_{jh}(x)=\frac1d(1-|4x_h-j|)_+,
 \qquad
 \phi^{\rm B}_{jh}(x)=\frac1d{4\choose j}x_h^j(1-x_h)^{4-j},
 \quad 0\le j\le4,
\]
with $1\le h\le d$. Both are nonnegative partitions of unity with $5d$ box readout coefficients. The first is an additive hat basis with an exact ReLU realization; the second is a degree-four Bernstein basis. Both purchase the same verified boundary contexts and approximate action labels, and optimize the same certificate loss. The label-loss control changes the objective, not its hat features. The achieved rational optimization certificate is evaluated after quantization.

\subsection{Complete recovery under a validation law}
The validation design draws each coordinate as $U/2^{16}$, with $U$ uniform on $\{0,\ldots,2^{16}-1\}$. These states are representable by the twenty-four-bit policy acquisition contract, but are not uniform over that finer grid. The following result applies to the declared independent context law and recovery. It does not assert that on-policy recursive states have that law.
\begin{proposition}[A finite catalogue of complete work]\label{prop:complete-work69}
Condition on a frozen finite family of predictors and a predictor-independent validation law. Suppose a failed screen invokes a deterministic recovery whose root-query count is at most $K(z)$, and every root query costs at most $W(z)$ under the declared descendant and arithmetic contract. If $0\le W(z)K(z)\le C_{\max}$, then with probability at least $1-\alpha$, simultaneously for $M$ fixed predictors,
\[
 \mathbb E[\text{recovery work}_\theta]
 \le\frac1n\sum_{i=1}^n W(z_i)K(z_i)\ell_\tau(z_i,\theta)
       +C_{\max}\sqrt{\frac{\log(M/\alpha)}{2n}}.
\]
Training, context acquisition, prediction, loading, identity checks and output costs are additional terms in complete work. Replacing $W$ by one gives a root-query result, not a bound on elapsed process time.
\end{proposition}
\noindent\emph{Proof.} See Supplement, Section~\ref{supp:retained69-proof1}.

The executed validation uses the root-query specialization and records descendant queries separately. It does not fit a runtime model and declare it a valid worst-case endpoint cost. The predictable-deployment result in Section~\ref{sec:deployment70} adds an explicit context-transfer premise and a complete operation account.

\subsection{An additional vector witness}
In the two-control economy a trained, projected and downward-rounded proposal can supply a new evaluated upper witness. Its insertion leaves the triangular lower cover unchanged. The original corner endpoints remain in that cover.
\begin{proposition}[Optional vector witnesses]\label{prop:vector-witness69}
Suppose a triangular action cover has a valid global lower bound $L$ and feasible evaluated witnesses with upper bound $U$. Adding any finite number of feasible, correctly evaluated witnesses and replacing $U$ by the minimum of all their upper bounds preserves containment of the original feasible infimum. Every returned witness with $U-L\le\delta$ has original Bellman regret at most $\delta$. The conclusion is independent of the training method and of whether an added witness is a vertex of the lower cover.
\end{proposition}
\noindent\emph{Proof.} See Supplement, Section~\ref{supp:retained69-proof2}.
The vector benchmark changes the trained proposal or the choice of a triangle to refine, not the authority of the global lower cover. Its training objective remains action-label loss.
